% Hanzo: A Crypto-Commerce Platform for Every Business
\documentclass[11pt,letterpaper]{article}
\usepackage{shared/hanzocover}
\usepackage{amsmath, amssymb}
\usepackage{graphicx}
\usepackage{booktabs}
\usepackage{hyperref}
\usepackage{enumitem}
\usepackage{xcolor}
\usepackage{float}
\hypersetup{colorlinks=true,linkcolor=black,citecolor=blue,urlcolor=blue}

\title{Hanzo: A Crypto-Commerce Platform for Every Business}
\author{
    Zach Kelling \\
    \textit{Hanzo AI (Techstars '17, Kansas City)} \\
    \texttt{zach@hanzo.ai}
}
\date{2018}

\begin{document}
\hanzocoverpage

\maketitle

\begin{abstract}
We describe Hanzo, a unified platform for cryptocurrency payments, smart
contract deployment, and digital wallet management targeting businesses with no
prior blockchain expertise.  The system exposes a single REST API through which
merchants accept Bitcoin and Ethereum payments with instant fiat settlement,
deploy parameterized smart contracts from an audited template library, and
manage both custodial and non-custodial wallets.  Two client-side JavaScript
libraries---Shop.js for checkout UI components and Coin.js for token
operations---reduce frontend integration to fewer than twenty lines of code.
Settlement occurs in the merchant's local currency within one business day at a
rate locked at transaction confirmation time, eliminating price-volatility risk
entirely.  The platform supports ERC-20 fungible tokens~\cite{vogelsteller2015}
and ERC-721 non-fungible tokens~\cite{entriken2018}, with an architecture
designed for extension to additional chains and token standards as the ecosystem
matures.
\end{abstract}

\section{Introduction}
\label{sec:intro}

By early 2018, global cryptocurrency market capitalization exceeded \$800
billion and the number of active wallet addresses surpassed 25 million.  Yet
merchant infrastructure lagged far behind consumer adoption.  The typical
integration path required a business to: (1) operate or interface with a
full blockchain node; (2) manage private keys and transaction signing;
(3) handle confirmation latency and chain reorganization risk; (4) convert
volatile cryptocurrency proceeds to fiat on an exchange; and (5) implement
regulatory compliance checks.  Each step introduced operational risk, and the
aggregate complexity deterred all but the most technically capable merchants.

Stripe demonstrated that the best payment integration is one where the
developer writes as little payments code as possible~\cite{stripe2018}.
Hanzo applies this principle to cryptocurrency.  A single API call creates a
payment intent; a single webhook confirms settlement.  Between those two
events, the platform handles address generation, confirmation monitoring,
exchange execution, fiat conversion, and bank transfer.

Bitcoin~\cite{nakamoto2008} proved that trustless value transfer is possible.
Ethereum~\cite{buterin2014} proved that programmable contracts can automate
business logic on-chain.  Hanzo provides the commercial infrastructure layer
that makes both accessible to any business, regardless of blockchain literacy.

\section{Architecture}
\label{sec:architecture}

The platform is organized into three subsystems: a payment gateway, a smart
contract engine, and an operations dashboard.  All three share a common
authentication layer, a unified data model, and a single set of API
credentials.

\subsection{Payment Gateway}

The payment gateway accepts cryptocurrency and settles in fiat.  Its
responsibilities are:

\begin{enumerate}[leftmargin=1.1em]
    \item \textbf{Address generation.}  For each payment intent, the gateway
    derives a unique deposit address from a hierarchical deterministic (HD)
    wallet~\cite{bip32}.  Address reuse is never permitted.
    \item \textbf{Rate locking.}  At intent creation time, the gateway queries
    partnered exchange order books and locks a conversion rate for a
    configurable window (default: 15 minutes).  The locked rate includes a
    spread that covers execution slippage and platform fees.
    \item \textbf{Confirmation monitoring.}  The gateway tracks the deposit
    address across the relevant blockchain and waits for a configurable number
    of confirmations (default: 1 for Ethereum, 3 for Bitcoin).  For
    high-value transactions, merchants may require additional confirmations.
    \item \textbf{Exchange execution.}  Upon sufficient confirmations, the
    gateway executes a market sell on one or more partnered exchanges.  The
    sell is atomic from the merchant's perspective: if the locked rate cannot
    be met, the platform absorbs the difference from its liquidity reserves.
    \item \textbf{Fiat settlement.}  Net proceeds are transferred to the
    merchant's bank account via ACH (US), SEPA (EU), or SWIFT (international),
    typically within one business day.
\end{enumerate}

\subsection{Smart Contract Engine}

The smart contract engine manages the lifecycle of on-chain programs.  It
stores a library of audited Solidity templates, each accepting a parameter
struct via the Hanzo API.  The engine compiles, deploys, and verifies the
contract on the target chain, then monitors its execution.  Merchants interact
with smart contracts entirely through REST endpoints; no Solidity toolchain
is required.

\subsection{Operations Dashboard}

The Hanzo Dashboard is a web application that provides real-time visibility
into all platform operations.  It displays payment status across currencies
and geographies, smart contract state and event logs, token sale progress with
investor-level detail, and aggregated business metrics.  The dashboard exposes
the same data available through the API, formatted for non-technical operators.

\section{Developer SDKs}
\label{sec:sdks}

The platform provides server-side libraries in Node.js, Python, Go, Ruby, and
Haskell, each wrapping the REST API with language-idiomatic types and error
handling.  Two client-side JavaScript libraries handle frontend integration.

\subsection{Shop.js}

Shop.js provides pre-built UI components for payment and checkout flows.
Components handle input validation, card number formatting, browser capability
detection, and responsive layout.  A complete checkout form requires importing
the library and mounting a single component:

\begin{verbatim}
import { Checkout } from '@hanzo/shop.js'

Checkout.mount('#payment-form', {
  apiKey: 'pk_live_...',
  currency: 'usd',
  amount: 4999
})
\end{verbatim}

The library renders a form that accepts both fiat card payments and
cryptocurrency, routing each to the appropriate backend flow.  Style
customization is handled through CSS custom properties.

\subsection{Coin.js}

Coin.js is a lightweight library for token-specific operations: accepting
cryptocurrency payments inline, displaying wallet balances, and interacting
with token sale contracts.  It wraps the Web3 provider detection and fallback
logic that would otherwise require significant boilerplate:

\begin{verbatim}
import { TokenSale } from '@hanzo/coin.js'

const sale = new TokenSale({
  apiKey: 'pk_live_...',
  contract: 'sale_abc123'
})
sale.mount('#token-purchase')
\end{verbatim}

Both libraries share a common event system, allowing developers to compose
fiat and crypto payment flows within a single checkout experience.

\section{Smart Contract Templates}
\label{sec:templates}

The template library covers four categories of on-chain business logic.  Each
template is deployed through the API by supplying a JSON parameter object; the
engine handles compilation, deployment, and source verification on Etherscan.

\subsection{Token Sales (ICOs)}

The ICO template implements a configurable crowdsale contract compliant with
the ERC-20 standard~\cite{vogelsteller2015}.  Parameters include:

\begin{itemize}[leftmargin=1.1em]
    \item Hard cap and soft cap denominated in ETH or USD-equivalent.
    \item Pricing tiers with time-based or volume-based transitions.
    \item Whitelist enforcement via Merkle root.
    \item Vesting schedules for team and advisor allocations.
    \item Refund mechanism if the soft cap is not met.
\end{itemize}

The contract is pausable by the contract owner and emits standard ERC-20
\texttt{Transfer} events for wallet compatibility.

\subsection{Airdrops}

The airdrop template distributes tokens to a list of addresses.  Rather than
executing individual transfers (which costs $O(n)$ gas for $n$ recipients),
the contract uses a Merkle-proof-based claiming mechanism.  The deployer
uploads a recipient list; the engine computes the Merkle root and stores it
on-chain.  Each recipient submits a proof to claim their allocation, paying
their own gas.  Total deployment cost is constant regardless of recipient
count.

\subsection{Collectibles}

The collectibles template deploys ERC-721 non-fungible token (NFT)
contracts~\cite{entriken2018}.  Each token carries a metadata URI pointing
to an off-chain JSON descriptor (name, description, image, attributes).  The
template supports:

\begin{itemize}[leftmargin=1.1em]
    \item Sequential or random minting.
    \item Royalty configuration for secondary sales.
    \item Marketplace listing hooks for integration with third-party platforms.
    \item Metadata freezing after a configurable period.
\end{itemize}

\subsection{Referral Programs}

The referral template implements on-chain referral tracking.  When a referred
user completes a qualifying action (purchase, registration, or token
acquisition), the contract automatically distributes a configurable reward to
the referrer's address.  Referral chains of arbitrary depth are supported, with
diminishing reward percentages at each level.

\section{Wallet Solutions}
\label{sec:wallets}

The platform supports two wallet modes, selectable per-user or per-transaction
through a single API parameter.

\subsection{Custodial Wallets}

In custodial mode, Hanzo manages private keys in a hardware security module
(HSM) cluster.  Key material never leaves the HSM boundary.  Transaction
signing occurs server-side, and wallet recovery uses standard identity
verification (email, phone, government ID).  This mode targets consumer-facing
applications where ease of use outweighs the trust assumption.

\subsection{Non-Custodial Wallets}

In non-custodial mode, keys are generated and stored on the user's device.
The Hanzo SDK provides key derivation (BIP-32/BIP-44), mnemonic backup
(BIP-39), and transaction construction, but never transmits private key
material to Hanzo servers.  Signing occurs client-side.  This mode targets
users who demand full sovereignty over their funds.

\begin{table}[H]
\centering
\small
\begin{tabular}{lcc}
\toprule
\textbf{Property} & \textbf{Custodial} & \textbf{Non-Custodial} \\
\midrule
Key storage & HSM & User device \\
Signing & Server-side & Client-side \\
Recovery & Identity verification & Mnemonic phrase \\
Trust model & Hanzo as custodian & Self-custody \\
API surface & Identical & Identical \\
\bottomrule
\end{tabular}
\caption{Wallet mode comparison.  Both modes expose the same API; the mode
parameter controls key management and signing behavior.}
\label{tab:wallets}
\end{table}

\section{Settlement and Foreign Exchange}
\label{sec:settlement}

Settlement is the operation that distinguishes Hanzo from a raw blockchain
interface.  The merchant receives fiat; Hanzo absorbs all cryptocurrency
exposure.

\subsection{Rate Locking}

At payment intent creation, the gateway samples order book depth across
partnered exchanges and computes a guaranteed rate that includes:

\begin{enumerate}[leftmargin=1.1em]
    \item The mid-market rate at the time of intent creation.
    \item A slippage buffer based on the payment amount and current liquidity.
    \item The platform fee (configurable per merchant, default 1\%).
\end{enumerate}

The locked rate is valid for the confirmation window.  If the market moves
adversely during confirmation, the platform's liquidity reserve absorbs the
difference.  If the market moves favorably, the surplus accrues to the
reserve.  Over a large volume of transactions, the reserve is
self-replenishing.

\subsection{FX Conversion}

For merchants settling in non-USD currencies, the gateway chains two
conversions: crypto-to-USD on the partnered exchange, then USD-to-local via
institutional FX rails.  The combined rate is quoted at intent creation and
locked for the confirmation window, so the merchant sees a single, stable
number in their local currency regardless of how many intermediate conversions
occur.

\subsection{Settlement Rails}

\begin{table}[H]
\centering
\small
\begin{tabular}{lll}
\toprule
\textbf{Region} & \textbf{Rail} & \textbf{Typical Latency} \\
\midrule
United States & ACH & T+1 \\
EU / EEA & SEPA & T+1 \\
United Kingdom & Faster Payments & Same day \\
International & SWIFT & T+2--3 \\
\bottomrule
\end{tabular}
\caption{Settlement rails by merchant region.}
\label{tab:settlement}
\end{table}

\section{Multi-Chain Support}
\label{sec:multichain}

At launch, the platform supports two chains: Bitcoin (UTXO model) and
Ethereum (account model, including all ERC-20 and ERC-721 tokens).  The
architecture is designed for chain-agnostic extension.

Each supported chain is represented by a \emph{chain adapter}---a module that
implements a standard interface:

\begin{itemize}[leftmargin=1.1em]
    \item \texttt{generateAddress(hdPath)} --- derive a deposit address.
    \item \texttt{watchAddress(addr, callback)} --- monitor for incoming
    transactions.
    \item \texttt{confirmTransaction(txHash, minConfirms)} --- wait for
    finality.
    \item \texttt{broadcastTransaction(signedTx)} --- submit a signed
    transaction.
\end{itemize}

Adding a new chain requires implementing this interface and registering the
adapter with the gateway.  The payment API, settlement logic, and merchant
dashboard are chain-agnostic; they operate on the abstract transaction model
produced by the adapter.

At the time of writing, adapter development is underway for Litecoin,
Bitcoin Cash, and EOS.  The ERC-20 adapter automatically supports any token
deployed on the Ethereum mainnet without per-token configuration.

\section{Deployment}
\label{sec:deployment}

Hanzo is a managed platform.  Merchants do not operate infrastructure.  The
production deployment consists of:

\begin{itemize}[leftmargin=1.1em]
    \item \textbf{API servers.}  Stateless HTTP services behind a load
    balancer, horizontally scaled by request volume.
    \item \textbf{Chain watchers.}  One process per supported chain, connected
    to geographically distributed full nodes for redundancy.  Watchers are
    single-writer to a shared transaction log, ensuring exactly-once
    processing.
    \item \textbf{Exchange connectors.}  Persistent WebSocket connections to
    partnered exchange APIs for real-time order book data and trade execution.
    \item \textbf{HSM cluster.}  Dedicated hardware security modules for
    custodial key storage, deployed in a quorum configuration requiring $k$
    of $n$ modules to sign.
    \item \textbf{Settlement engine.}  A batch processor that aggregates
    confirmed transactions, computes net settlement amounts per merchant, and
    initiates bank transfers on a daily cycle.
    \item \textbf{Dashboard.}  A single-page application served from a CDN,
    consuming the same API available to developers.
\end{itemize}

All components emit structured logs and metrics to a centralized observability
stack.  Uptime is monitored at \texttt{status.hanzo.ai} with public incident
reporting.

\subsection{API Versioning}

The API is versioned with a date-based scheme.  Breaking changes are released
under a new version header; older versions continue to function for a minimum
of twelve months.  Non-breaking additions (new fields, new endpoints) are
available immediately on all versions.  This policy ensures that merchants
upgrade on their own schedule without surprises.

\section{Security}
\label{sec:security}

\begin{itemize}[leftmargin=1.1em]
    \item \textbf{Smart contract audits.}  Every template undergoes
    third-party audit before inclusion in the library.  Audit reports are
    published and linked from the API documentation.
    \item \textbf{Fraud detection.}  Transaction pattern analysis flags
    anomalous behavior (unusual amounts, rapid-fire transactions, known
    tainted addresses) and holds the transaction for manual review.
    \item \textbf{Key security.}  Custodial keys are stored exclusively in
    HSMs (FIPS 140-2 Level 3).  Non-custodial keys never leave the client
    device.
    \item \textbf{Transport security.}  All API traffic is TLS~1.2+ with
    certificate pinning available in mobile SDKs.
    \item \textbf{Compliance.}  KYC/AML checks are integrated at the
    account level, with jurisdiction-specific rule sets configurable per
    merchant.
\end{itemize}

\section{Conclusion}
\label{sec:conclusion}

Hanzo reduces cryptocurrency commerce to an API integration.  A merchant adds
a few lines of JavaScript, receives fiat in their bank account, and gains
access to smart contract-powered business models---token sales, airdrops,
collectibles, referral programs---without writing Solidity or operating
blockchain infrastructure.  The zero-volatility settlement model eliminates
the primary financial risk of accepting cryptocurrency, and the dual
custodial/non-custodial wallet architecture lets businesses serve both
mainstream consumers and self-sovereign users through a single integration.

The platform is operational and accepting merchants.  Documentation is
available at \texttt{docs.hanzo.ai}; developer support is available at
\texttt{\#hanzo} on Freenode.

\begin{thebibliography}{9}

\bibitem{nakamoto2008}
S.~Nakamoto.
\newblock Bitcoin: A peer-to-peer electronic cash system.
\newblock 2008.

\bibitem{buterin2014}
V.~Buterin.
\newblock Ethereum: A next-generation smart contract and decentralized
application platform.
\newblock 2014.

\bibitem{vogelsteller2015}
F.~Vogelsteller and V.~Buterin.
\newblock ERC-20 token standard.
\newblock Ethereum Improvement Proposal 20, 2015.

\bibitem{entriken2018}
W.~Entriken, D.~Shirley, J.~Evans, and N.~Sachs.
\newblock ERC-721 non-fungible token standard.
\newblock Ethereum Improvement Proposal 721, 2018.

\bibitem{bip32}
P.~Wuille.
\newblock BIP-32: Hierarchical deterministic wallets.
\newblock Bitcoin Improvement Proposal 32, 2012.

\bibitem{stripe2018}
Stripe, Inc.
\newblock Stripe payments documentation.
\newblock \url{https://stripe.com/docs}, 2018.

\end{thebibliography}

\end{document}
